% ============================================
% MAZE SOLVER ROBOT - PROJECT PROPOSAL
% ============================================
\documentclass[aspectratio=169,11pt]{beamer}

% --- Packages ---
\usepackage[utf8]{inputenc}
\usepackage[T1]{fontenc}
\usepackage{lmodern}
\usepackage{amsmath,amssymb}
\usepackage{graphicx}
\usepackage{tikz}
\usetikzlibrary{shapes.geometric, arrows.meta, positioning, calc, fit, backgrounds, shadows}
\usepackage{circuitikz}
\usepackage{booktabs}
\usepackage{array}
\usepackage{multirow}
\usepackage{xcolor}
\usepackage{colortbl}  % REQUIRED for \cellcolor
\usepackage{listings}
\usepackage{hyperref}

% --- Theme & Colors ---
\usetheme{moloch}


% For the title page specifically, also set these:
\setbeamercolor{title}{fg=teal!80!black}
\setbeamercolor{subtitle}{fg=teal!60!black}
\setbeamercolor{author}{fg=black}
\setbeamercolor{institute}{fg=gray}
\setbeamercolor{date}{fg=gray}


% --- Title Info ---
\title[Maze Solver Robot]{Maze Solver Robot}
\subtitle{Project Proposal --- Autonomous Maze Exploration and Shortest-Path Solving on ATmega32}
\author{Md. Ariful Islam \and Hasan Ferdous Mahee \and B. M. Sufian Ashraf}
\institute{Department of Computer Science and Engineering}
\date{\today}

% --- Custom Commands ---
\newcommand{\highlight}[1]{\textcolor{teal!80!black}{\textbf{#1}}}
\newcommand{\pin}[1]{\texttt{\textcolor{orange!80!black}{#1}}}

% ============================================
\begin{document}
% ============================================

% --- Title Slide ---
\begin{frame}
    \titlepage
\end{frame}

% --- Outline ---
\begin{frame}{Outline}
    \tableofcontents
\end{frame}

% ============================================
\section{1. Motivation}
% ============================================
\begin{frame}{Problem Statement}
    \small
    \begin{block}{What is the Problem?}
        \small Autonomous navigation through an \highlight{unknown environment} is a core problem in mobile robotics --- search-and-rescue robots, warehouse pickers, and cleaning robots all need to explore, build an internal map, and find an efficient path without any prior knowledge of the layout or human guidance.
    \end{block}

    \vspace{0.25cm}
    \begin{columns}[T]
        \begin{column}{0.48\textwidth}
            \textbf{Why is it Important?}
            \begin{itemize}\setlength{\itemsep}{0pt}
                \item Core capability behind search-and-rescue robotics
                \item Foundation of SLAM and exploration algorithms
                \item Directly transfers to warehouse/indoor navigation
            \end{itemize}
        \end{column}
        \begin{column}{0.48\textwidth}
            \textbf{Why This Project?}
            \begin{itemize}\setlength{\itemsep}{0pt}
                \item Combines real-time control with classical search algorithms
                \item Forces on-the-fly mapping under memory constraints
                \item Natural, well-scoped extension of embedded control (ATmega32)
            \end{itemize}
        \end{column}
    \end{columns}
\end{frame}

\begin{frame}{Project Objectives}
    \begin{block}{Main Objectives}
        \begin{enumerate}
            \item \highlight{Design and build} an autonomous robot capable of exploring an unknown maze of unit cells using wall-following sensors.
            \item \highlight{Implement wall detection} using ultrasonic sensors on the front, left, and right of the robot.
            \item \highlight{Implement an exploration algorithm} (flood-fill) that incrementally maps the maze and updates the shortest known path to the goal.
            \item \highlight{Implement a run-time architecture:} explore the maze on the first run, then execute the fastest known path on a speed run.
            \item \highlight{Develop precise motor control} with PWM speed regulation and closed-loop odometry via optical wheel encoders for accurate cell-to-cell moves and turns.
            \item \highlight{Handle the dead-end edge case:} detect dead ends during exploration and backtrack correctly to continue the search.
        \end{enumerate}
    \end{block}
\end{frame}

% ============================================
\section{2. Equipment}
% ============================================
\begin{frame}{Hardware Components}
    \begin{table}
        \centering
        \tiny
        \begin{tabular}{@{}lp{7.4cm}@{}}
            \toprule
            \textbf{Component} & \textbf{Purpose} \\
            \midrule
            ATmega32A & Central microcontroller; runs flood-fill, wall sensing, and motor control \\
            3x HC-SR04 Ultrasonic Sensors & Front, left, and right wall detection at each cell \\
            2x Wheel Encoders (Optical) & Optical encoder discs + sensors on each wheel for precise distance and turn odometry \\
            L298N Motor Driver Module & Dual H-bridge module for PWM motor speed and direction control \\
            DC TT Gear Motors & Provides differential drive locomotion \\
            2WD Chassis + Wheels & Mechanical frame, wheels, and caster for mobility and stability \\
            Status LEDs & Visual indication of robot state (exploring, dead-end, goal reached, speed run) \\
            L7805 Voltage Regulator & Stabilizes 5V supply for MCU, sensors, and logic \\
            3.7V 18650 Li-Ion Batteries (x3) & 3S configuration providing ~11.1V main power supply \\
            18650 Battery Holder (3S) & Houses batteries and supplies power to driver and regulator \\
            Crystal Oscillator (16 MHz) & External clock source for ATmega32A system clock \\
            AVR ISP Programmer & In-System Programmer (ISP) for flashing code onto ATmega32A \\
            \bottomrule
        \end{tabular}
    \end{table}
\end{frame}

\begin{frame}{Development Tools}
    \scriptsize
    \textbf{Software and Programming Tools}
    \begin{itemize}\setlength{\itemsep}{0pt}
        \item \highlight{Atmel Studio 7.0 / AVR-GCC} --- Writing and compiling the embedded C program
        \item \highlight{AVRDUDE} --- Uploading the compiled firmware to the ATmega32
        \item \highlight{Git} --- Version control and collaboration for the source code
    \end{itemize}

    \vspace{0.12cm}
    \textbf{Hardware Development Tools}
    \begin{itemize}\setlength{\itemsep}{0pt}
        \item \highlight{AVR ISP Programmer} --- Programming the ATmega32 through ISP
        \item \highlight{Breadboard + Jumper Wires} --- Circuit assembly, component connection, and hardware testing
    \end{itemize}

    \vspace{0.12cm}
    \begin{alertblock}{Note}
        \scriptsize All development will be done in C using AVR Libc, without the Arduino framework, allowing direct access to ATmega32 hardware registers --- including Timer1 input capture for reading three ultrasonic echoes.
    \end{alertblock}
\end{frame}

% ============================================
\section{3. System Design}
% ============================================
\begin{frame}{System Block Diagram}
    \centering

    \definecolor{themeTeal}{RGB}{8,92,92}
    \definecolor{powerFill}{RGB}{235,246,240}
    \definecolor{sensorFill}{RGB}{231,243,246}
    \definecolor{mcuFill}{RGB}{244,235,241}
    \definecolor{driverFill}{RGB}{248,241,230}
    \definecolor{motorFill}{RGB}{249,237,237}
    \definecolor{sonarFill}{RGB}{255,236,236}
    \definecolor{indicatorFill}{RGB}{255,244,224}

    \begin{tikzpicture}[
        box/.style={
            rectangle,
            rounded corners=5pt,
            draw=themeTeal,
            line width=0.9pt,
            minimum width=1.9cm,
            minimum height=0.68cm,
            align=center,
            font=\tiny\bfseries,
            drop shadow={
                opacity=0.15,
                shadow xshift=1pt,
                shadow yshift=-1pt
            }
        },
        sonarbox/.style={
            box,
            fill=sonarFill,
            minimum height=0.56cm
        },
        line/.style={
            line width=1.1pt,
            draw=themeTeal!85!black
        },
        bigbox/.style={
            rectangle,
            rounded corners=8pt,
            draw=themeTeal,
            line width=1.2pt,
            minimum width=2.6cm,
            minimum height=2.5cm,
            align=center,
            font=\small\bfseries,
            fill=mcuFill,
            drop shadow={
                opacity=0.18,
                shadow xshift=1.5pt,
                shadow yshift=-1.5pt
            }
        },
        arrow/.style={
            -{Stealth[length=2.2mm]},
            line width=1pt,
            draw=themeTeal
        },
        powerarrow/.style={
            -{Stealth[length=2.2mm]},
            line width=1.1pt,
            draw=themeTeal!85!black
        },
        label/.style={
            font=\tiny\bfseries,
            fill=white,
            draw=themeTeal!45,
            rounded corners=2pt,
            inner sep=1.2pt
        }
    ]

        % Power supply
        \node[box, fill=powerFill] at (-5.6,2.5) (battery)
            {3x 3.7V Li-Ion\\($\sim$11.1V)};

        \node[box, fill=powerFill] at (-3.3,2.5) (reg)
            {L7805\\5V Reg.};

        \draw[powerarrow] (battery.east) -- (reg.west);

        % Microcontroller
        \node[bigbox] at (0,0) (mcu)
            {\textbf{ATmega32A}\\[3pt]
             \scriptsize GPIO $\mid$ PWM\\[2pt]
             \scriptsize Timer1 Input Capture\\[2pt]
             \tiny (16MHz Clock)};

        % Regulated 5V connection
        \coordinate (vfive) at (0,1.4);

        \draw[line]
            (reg.south) -- ++(0,-0.3) -- (vfive)
            node[pos=0.75, label] {+5V};

        \draw[powerarrow]
            (vfive) -- (mcu.north);

        % Three ultrasonic sensors: front, left, right
        \node[sonarbox, minimum width=2.1cm] at (-5.4, 1.35) (sonarL)
            {HC-SR04\\Left Wall};

        \node[sonarbox, minimum width=2.1cm] at (0, 2.55) (sonarF)
            {HC-SR04\\Front Wall};

        \node[sonarbox, minimum width=2.1cm] at (-5.4, -1.35) (sonarR)
            {HC-SR04\\Right Wall};

        \draw[arrow] (sonarL.east) -- ++(0.9,0)
            |- ([yshift=0.55cm]mcu.west)
            node[pos=0.32, above, label] {TRIG/ECHO};

        \draw[arrow] (sonarF.south) -- ([xshift=-0.3cm]mcu.north)
            node[pos=0.5, right=1pt, label] {TRIG/ECHO};

        \draw[arrow] (sonarR.east) -- ++(0.9,0)
            |- ([yshift=-0.55cm]mcu.west)
            node[pos=0.32, below, label] {TRIG/ECHO};

        % Status LED indicator block
        \node[box, fill=indicatorFill, minimum width=2.1cm] at (0,-2.55) (ind)
            {Status LEDs};

        \draw[arrow] (mcu.south) -- ([xshift=-0.3cm]ind.north)
            node[pos=0.5, right=1pt, label] {GPIO};

        % Motor driver
        \node[
            box,
            fill=driverFill,
            minimum width=2.1cm,
            minimum height=0.9cm
        ] at (4.3,0) (driver)
            {L298N\\Motor Driver};

        \draw[arrow] (mcu.east) -- (driver.west)
            node[midway, above=2pt, label] {PWM + DIR};

        % Motor-supply route
        \draw[powerarrow]
            (battery.north)
            -- ++(0,0.4)
            -| (driver.north)
            node[pos=0.68, above, label] {Motor Supply (11.1V)};

        % Motors
        \node[box, fill=motorFill] at (7.1,0.68) (motL)
            {Left DC\\TT Motor};

        \node[box, fill=motorFill] at (7.1,-0.68) (motR)
            {Right DC\\TT Motor};

        \draw[arrow]
            ([yshift=0.17cm]driver.east)
            -- ++(0.32,0)
            |- (motL.west);

        \draw[arrow]
            ([yshift=-0.17cm]driver.east)
            -- ++(0.32,0)
            |- (motR.west);

        % Wheel encoders (optical), mounted alongside the wheels
        \node[box, fill=sonarFill, minimum width=1.7cm, minimum height=0.5cm] at (7.1,1.55) (encL)
            {Left Encoder};

        \node[box, fill=sonarFill, minimum width=1.7cm, minimum height=0.5cm] at (7.1,-1.55) (encR)
            {Right Encoder};

        \draw[arrow] (motL.north) -- (encL.south);
        \draw[arrow] (motR.south) -- (encR.north);

        \draw[arrow] (encL.west) -| ([xshift=-0.2cm]mcu.north)
            node[pos=0.25, above, label] {\tiny Pulses};

        \draw[arrow] (encR.west) -| ([xshift=-0.2cm]ind.south)
            node[pos=0.25, below, label] {\tiny Pulses};

    \end{tikzpicture}
\end{frame}

\begin{frame}[t]{Component Roles}
    \scriptsize
    \begin{columns}[T]
        \begin{column}{0.48\textwidth}
            \begin{block}{3x HC-SR04 Ultrasonic Sensors}
                \begin{itemize}\setlength{\itemsep}{0pt}\setlength{\parskip}{0pt}
                    \item Mounted front, left, right; each returns distance-to-wall
                    \item TRIG pulse $\rightarrow$ ECHO pulse width $\propto$ distance
                    \item A cell wall is flagged present below a calibrated fixed threshold distance
                    \item Timer1 input capture multiplexed across the three sensors for precise timing
                \end{itemize}
            \end{block}

            \begin{block}{L298N Motor Driver Module}
                \begin{itemize}\setlength{\itemsep}{0pt}\setlength{\parskip}{0pt}
                    \item Dual H-bridge module (drives 2 motors)
                    \item Direction control via IN1--IN4 logic pins
                    \item Speed control via ATmega32A PWM on ENA/ENB
                \end{itemize}
            \end{block}
        \end{column}

        \begin{column}{0.48\textwidth}
            \begin{block}{DC TT Gear Motors + Optical Encoders}
                \begin{itemize}\setlength{\itemsep}{0pt}\setlength{\parskip}{0pt}
                    \item 3V--6V rated motors (~150--200 RPM)
                    \item Each wheel paired with an optical encoder for pulse counting
                    \item Differential drive: 90\textdegree{} turns and cell-to-cell moves measured by encoder pulse count, not timing
                \end{itemize}
            \end{block}

            \begin{block}{Status LEDs}
                \begin{itemize}\setlength{\itemsep}{0pt}\setlength{\parskip}{0pt}
                    \item Indicate current robot state: exploring, dead-end/backtrack, goal reached, speed run
                \end{itemize}
            \end{block}

            \begin{block}{Power System}
                \begin{itemize}\setlength{\itemsep}{0pt}\setlength{\parskip}{0pt}
                    \item 3x 3.7V Li-Ion (3S): ~11.1V pack; L7805 steps down to clean 5V logic rail
                \end{itemize}
            \end{block}
        \end{column}
    \end{columns}
\end{frame}

\begin{frame}{ATmega32 Pin Assignment}
    \begin{table}
        \centering
        \tiny
        \begin{tabular}{@{}lllp{4.6cm}@{}}
            \toprule
            \textbf{Function} & \textbf{Port} & \textbf{Pins} & \textbf{Description} \\
            \midrule
            Front Sonar & PORTB & PB0 (TRIG), PB1/ICP1 (ECHO) & HC-SR04 front wall ranging \\
            Left Sonar & PORTA & PA5 (TRIG), PA6 (ECHO) & HC-SR04 left wall ranging \\
            Right Sonar & PORTA & PA7 (TRIG), PC7 (ECHO) & HC-SR04 right wall ranging \\
            Motor Direction L & PORTD, PORTC & PD7, PC0 & Left motor direction control (L298N IN1, IN2) \\
            Motor Direction R & PORTC & PC1, PC2 & Right motor direction control (L298N IN3, IN4) \\
            PWM Output L & PORTD & PD4 (OC1B) & Timer1 PWM speed control for left motor (ENA) \\
            PWM Output R & PORTD & PD5 (OC1A) & Timer1 PWM speed control for right motor (ENB) \\
            Left Encoder & PORTD & PD2 (INT0) & Left wheel optical encoder pulse input \\
            Right Encoder & PORTD & PD3 (INT1) & Right wheel optical encoder pulse input \\
            Status LEDs & PORTC & PC3--PC5 & Digital outputs for robot-state indicator LEDs \\
            ISP Programming & PORTB & PB5--PB7 & MOSI, MISO, SCK interface for AVR ISP \\
            Reset Circuit & --- & Pin 9 (RESET) & Active-low reset with 10k$\Omega$ pull-up resistor \\
            Clock Source & --- & Pins 12, 13 & XTAL1/XTAL2 with 16 MHz crystal \& 22pF caps \\
            \bottomrule
        \end{tabular}
    \end{table}
    \vspace{0.1cm}
    \scriptsize\textcolor{gray}{Note: only PB1 uses hardware input capture (ICP1) for the front sonar; left/right echoes are timed via polling. Left/right encoders use the two external interrupt pins (INT0, INT1) for reliable pulse counting.}
\end{frame}


% ============================================
\section{4. Role of the ATmega32}
% ============================================
\begin{frame}{ATmega32 --- Central Processing Unit}
    \small
    The ATmega32 is the \highlight{sole controller} of this project. It performs all sensing, mapping, decision-making, and actuation tasks --- there is no secondary microcontroller.

    \vspace{0.2cm}
    \begin{block}{Core Tasks Performed by ATmega32}
        \small
        \begin{enumerate}\setlength{\itemsep}{0pt}
            \item \textbf{Wall Sensing:} Triggers and times all three ultrasonic sensors each cell
            \item \textbf{Cell Mapping:} Records discovered walls into an in-memory maze grid
            \item \textbf{Flood-Fill Computation:} Recomputes distance-to-goal values as new walls are discovered
            \item \textbf{Motion Planning:} Chooses the next cell to move into based on current flood values
            \item \textbf{Motor Control:} Executes forward moves and in-place turns using encoder pulse counts for closed-loop accuracy
            \item \textbf{Run Management:} Tracks exploration vs. speed-run state and dead-end backtracking
        \end{enumerate}
    \end{block}
\end{frame}

\begin{frame}{ATmega32 Peripheral Usage}
    \scriptsize
    \setlength{\itemsep}{-2pt}
    \begin{columns}[T]
        \begin{column}{0.48\textwidth}
            \begin{block}{Digital I/O (Sonar Trigger/Echo)}
                \begin{itemize}\setlength{\itemsep}{-2pt}
                    \item \textbf{PORTA/PORTB/PORTC:} TRIG outputs and ECHO inputs for 3 sensors
                    \item Sensors triggered sequentially each cell to avoid cross-talk
                \end{itemize}
            \end{block}

            \vspace{-0.15cm}
            \begin{block}{Timer1 --- PWM \& Input Capture}
                \begin{itemize}\setlength{\itemsep}{-2pt}
                    \item \pin{OC1B} (PD4) $\rightarrow$ Left motor PWM speed (ENA)
                    \item \pin{OC1A} (PD5) $\rightarrow$ Right motor PWM speed (ENB)
                    \item \pin{ICP1} (PB1) $\rightarrow$ \highlight{Input capture for front-sonar echo timing}
                \end{itemize}
            \end{block}

            \vspace{-0.15cm}
            \begin{block}{GPIO (Motor Control)}
                \begin{itemize}\setlength{\itemsep}{-2pt}
                    \item PD7, PC0 $\rightarrow$ Left motor H-Bridge (IN1, IN2)
                    \item PC1, PC2 $\rightarrow$ Right motor H-Bridge (IN3, IN4)
                \end{itemize}
            \end{block}
        \end{column}
        \begin{column}{0.48\textwidth}
            \begin{block}{External Interrupts (Wheel Encoders)}
                \begin{itemize}\setlength{\itemsep}{-2pt}
                    \item \pin{INT0} (PD2) $\rightarrow$ Left wheel encoder pulse count
                    \item \pin{INT1} (PD3) $\rightarrow$ Right wheel encoder pulse count
                    \item Pulse counts drive closed-loop distance and turn control
                \end{itemize}
            \end{block}

            \vspace{-0.15cm}
            \begin{block}{SRAM --- Maze State}
                \begin{itemize}\setlength{\itemsep}{-2pt}
                    \item Wall bitmap: 4 bits/cell (N/E/S/W) for up to N$\times$N cells
                    \item Flood-fill distance array: 1 byte/cell
                    \item Fits comfortably within the ATmega32's 2KB SRAM
                \end{itemize}
            \end{block}

            \vspace{-0.15cm}
            \begin{block}{GPIO (Status LEDs)}
                \begin{itemize}\setlength{\itemsep}{-2pt}
                    \item \textbf{PC3--PC5:} indicate exploring, dead-end, and goal/speed-run states
                \end{itemize}
            \end{block}

            \vspace{-0.15cm}
            \begin{block}{Other Features}
                \scriptsize
                \textbf{Watchdog Timer:} reset on software hang
            \end{block}
        \end{column}
    \end{columns}
\end{frame}

% ============================================
\section{5. Maze Mapping and Algorithm}
% ============================================
\begin{frame}{Maze Representation and Flood-Fill}
    \scriptsize
    \begin{block}{Cell and Wall Representation}
        \begin{itemize}\setlength{\itemsep}{0pt}\setlength{\parskip}{0pt}
            \item The maze is modeled as an \highlight{$N \times N$ grid of cells}, each with up to 4 walls (N, E, S, W)
            \item Walls are discovered incrementally as the robot visits each cell and reads its 3 ultrasonic sensors
            \item Each cell's wall state is stored as a \highlight{4-bit bitmap}; the full maze fits in well under 1KB of SRAM
        \end{itemize}
    \end{block}

    \vspace{0.15cm}
    \begin{block}{Why Flood-Fill?}
        \begin{itemize}\setlength{\itemsep}{0pt}\setlength{\parskip}{0pt}
            \item Each cell holds an integer \highlight{distance-to-goal} value, initialized by Manhattan distance and updated by a BFS-style wave from the goal outward
            \item When a new wall is discovered, only the affected region is recomputed --- \highlight{cheap enough to rerun after every move} on an 8-bit MCU
            \item Guarantees the robot always moves toward a cell on a currently-shortest known path, without needing global replanning from scratch
        \end{itemize}
    \end{block}

    \vspace{0.15cm}
    \textbf{Operational Cycle:} \highlight{Run 1 (Exploration):} start from the entry cell, sense walls with the 3 sonars, explore the maze cell by cell, back out of dead ends when hit, and continue until the goal is reached --- the map and the shortest known path are remembered in SRAM. \highlight{Run 2 (Speed Run):} starting again from the entry cell, follow the remembered shortest path directly to the goal, with no further exploration.
\end{frame}

% ============================================
\section{6. Implementation Plan}
% ============================================

\begin{frame}{Expected Outcome}
    \begin{block}{Final Deliverable}
        A fully assembled autonomous maze-solving robot that:
        \begin{itemize}
            \item Explores an \highlight{unknown maze of unit cells} using front/left/right wall sensing
            \item Builds an internal map and computes the \highlight{shortest known path via flood-fill} as it explores
            \item Reaches the goal cell on its exploration run, then executes a \highlight{fast, direct speed run} along the shortest discovered path
            \item Reliably completes a maze of at least \highlight{3$\times$3 cells} within a bounded time and move budget
        \end{itemize}
    \end{block}
\end{frame}

\begin{frame}{Development Stages}
    \scriptsize
    \begin{columns}[T]
        \begin{column}{0.48\textwidth}
            \begin{block}{Stage 1: Project Update 1}
                \textbf{Milestones:}
                \begin{itemize}\setlength{\itemsep}{0pt}
                    \item[$\square$] Assemble chassis, motors, and control circuit
                    \item[$\square$] Integrate and calibrate 3x HC-SR04 sensors
                    \item[$\square$] Implement encoder-based forward-move and 90\textdegree{}-turn primitives
                    \item[$\square$] Sonars detect walls and show status via LEDs
                \end{itemize}
                \textbf{Deliverable:} A robot that reliably moves cell-to-cell and detects walls on three sides --- the base platform for maze solving
            \end{block}
        \end{column}

        \begin{column}{0.48\textwidth}
            \begin{block}{Stage 2: Final}
                \textbf{Milestones:}
                \begin{itemize}\setlength{\itemsep}{0pt}
                    \item[$\square$] Implement maze grid storage and wall-bitmap updates
                    \item[$\square$] Implement flood-fill computation and incremental updates
                    \item[$\square$] Implement full exploration run to goal with dead-end handling
                    \item[$\square$] Implement speed run along the shortest discovered path
                    \item[$\square$] Full-maze testing, calibration, and LED signaling
                \end{itemize}
                \textbf{Deliverable:} Complete maze-solver robot demonstrating exploration, mapping, and a faster speed run
            \end{block}
        \end{column}
    \end{columns}
\end{frame}

\begin{frame}{Group Member Responsibilities}
    \begin{table}
        \centering
        \footnotesize
        \begin{tabular}{@{}lp{7.5cm}@{}}
            \toprule
            \textbf{Member} & \textbf{Primary Responsibilities} \\
            \midrule
            Hasan Ferdous Mahee & Hardware assembly, power system, and sonar integration \\
            Md. Ariful Islam & Motion primitives (turns/moves) and wall-detection calibration \\
            B. M. Sufian Ashraf & Maze representation, flood-fill algorithm, and run-state management \\
            All Members & Testing, documentation, and presentation \\
            \bottomrule
        \end{tabular}
    \end{table}
\end{frame}

% ============================================
\section{7. Use of Other Controllers}
% ============================================
\begin{frame}{No Additional Controllers Used}
    \begin{block}{Design Decision}
        This project uses \highlight{only the ATmega32} microcontroller. No ESP32, Arduino, or Raspberry Pi is required.
    \end{block}

    \vspace{0.3cm}
    \textbf{Justification:}
    \begin{itemize}
        \item Flood-fill on a small grid is cheap enough in both memory and compute for an 8-bit MCU
        \item Sensor and motor I/O needs (3 sonar pairs, 2 PWM channels, 4 direction pins) fit well within the ATmega32's 32 available pins
        \item A single controller keeps the system simpler, cheaper, and easier to debug than a dual-controller design
        \item Direct register-level programming gives full control over sonar timing precision, which matters for reliable wall detection
    \end{itemize}
\end{frame}

% ============================================
\section{8. Testing and Success Criteria}
% ============================================
\begin{frame}{Testing Methodology}
    \small
    \begin{block}{Test Environment Setup}
        \small
        \begin{itemize}\setlength{\itemsep}{0pt}
            \item \textbf{Maze:} Physical grid maze with uniform unit cells and vertical walls for reliable sonar reflection
            \item \textbf{Maze Size:} 3$\times$3 cells, including one dead end and at least one intersection
        \end{itemize}
    \end{block}

    \vspace{0.1cm}
    \textbf{Test Procedures:}
    \begin{enumerate}\setlength{\itemsep}{0pt}
        \item \textbf{Unit Testing:} Verify each of the 3 sonar sensors independently against known distances
        \item \textbf{Module Testing:} Test encoder-based forward-move and turn accuracy over repeated cycles
        \item \textbf{Integration Testing:} Sonars detect walls correctly and status LEDs reflect robot state
        \item \textbf{Full-Maze Testing:} Complete exploration run, followed by a speed run; measure both times
        \item \textbf{Repeatability:} Run 5 consecutive full explore-then-speed-run cycles; record success/failure
    \end{enumerate}
\end{frame}

\begin{frame}{Success Criteria}
    \vspace{0.3cm}
    \begin{block}{Essential Features (Must Have)}
        \begin{itemize}\setlength{\itemsep}{4pt}
            \item[$\checkmark$] Correctly detects walls on 3 sides at each cell
            \item[$\checkmark$] Correctly detects and backtracks out of dead ends during exploration
            \item[$\checkmark$] Completes full exploration of a 3$\times$3 maze and reaches the goal
            \item[$\checkmark$] Flood-fill map stays consistent as walls are discovered
            \item[$\checkmark$] Executes a speed run strictly faster than the exploration run, following the remembered path directly
        \end{itemize}
    \end{block}
\end{frame}

% ============================================
\section{9. Challenges and Future Scope}
% ============================================

\begin{frame}{Expected Challenges}
    \begin{itemize}
        \item \highlight{Sonar cross-talk} between the three simultaneously-mounted ultrasonic sensors
        \item \highlight{Encoder calibration} for consistent pulse-count-to-distance and pulse-count-to-turn-angle mapping
        \item \highlight{Wall-detection thresholding} under sensor noise near cell boundaries
        \item \highlight{SRAM budgeting} for the wall bitmap and flood-fill arrays as maze size grows
        \item \highlight{Dead-end handling} in the exploration algorithm --- reliable detection and backtracking
    \end{itemize}
\end{frame}

\begin{frame}{Future Scope}
    \small
    \begin{columns}[T]
        \begin{column}{0.48\textwidth}
            \textbf{Short-Term Improvements}
            \begin{itemize}\setlength{\itemsep}{0pt}
                \item Integrate the \highlight{optical wheel encoders} (already acquired) for closed-loop odometry
                \item Add a \highlight{gyroscope} for precise turn correction
            \end{itemize}
        \end{column}
        \begin{column}{0.48\textwidth}
            \textbf{Long-Term Extensions}
            \begin{itemize}\setlength{\itemsep}{0pt}
                \item Full \highlight{micromouse-style} speed optimization (smooth curves, ramped acceleration)
                \item \highlight{Multi-robot maze solving} with shared map data
                \item Port to \highlight{dynamic/reconfigurable mazes}
                \item \highlight{Machine-learned} move policies as an alternative to flood-fill
            \end{itemize}
        \end{column}
    \end{columns}

    \vspace{0.2cm}
    \begin{alertblock}{Real-World Applications}
        Search-and-rescue robots, autonomous warehouse navigation, indoor drones, and educational micromouse competitions.
    \end{alertblock}
\end{frame}

% --- Thank You ---
\begin{frame}
    \begin{center}
        \Huge\textbf{Thank You!}

        \vspace{1.5cm}
        \normalsize
        \textit{Maze Solver Robot --- ATmega32 Project Proposal}\\[5pt]
        \textbf{Md. Ariful Islam \quad Hasan Ferdous Mahee \quad B. M. Sufian Ashraf}
    \end{center}
\end{frame}

% ============================================
\end{document}
